
\documentclass{article}
\usepackage{colortbl}
\usepackage{makecell}
\usepackage{multirow}
\usepackage{supertabular}

\begin{document}

\newcounter{utterance}

\centering \large Interaction Transcript for game `cladder', experiment `full\_v1.5\_default', episode 9887 with Qwen3.5{-}9B.
\vspace{24pt}

{ \footnotesize  \setcounter{utterance}{1}
\setlength{\tabcolsep}{0pt}
\begin{supertabular}{c@{$\;$}|p{.15\linewidth}@{}p{.15\linewidth}p{.15\linewidth}p{.15\linewidth}p{.15\linewidth}p{.15\linewidth}}
    \# & \multicolumn{2}{c}{Player} && \multicolumn{2}{c}{Game Master} \\
    \hline

    \theutterance \stepcounter{utterance}  
    & & & \multicolumn{4}{p{0.6\linewidth}}{
        \cellcolor[rgb]{0.9,0.9,0.9}{
            \makecell[{{p{\linewidth}}}]{
                \texttt{\tiny{[P1$\langle$GM]}}
                \texttt{You are an expert in causal inference. The following question is not a typical commonsense query, but rather a meticulously designed question created by a professor specializing in causal inference, intended to assess the students' mastery of the course content.} \\
\texttt{Imagine a self{-}contained, hypothetical world with only the following conditions, and without any unmentioned factors or causal relationships: Yupt has a direct effect on muvq. Zupj has a direct effect on muvq. We know that yupt or zupj causes muvq. We observed an individual is not zupj. Would an individual is muvq if yupt instead of not yupt?} \\
\texttt{Start your answer with 'Yes' or 'No', followed by additional reasoning or evidence to support your explanation.} \\
            }
        }
    }
     \\ \\

    \theutterance \stepcounter{utterance}  
    & \multicolumn{4}{p{0.6\linewidth}}{
        \cellcolor[rgb]{0.561,0.737,0.561}{
            \makecell[{{p{\linewidth}}}]{
                \texttt{\tiny{[P1$\rangle$GM]}}
                \texttt{No} \\
\\ 
\texttt{**Reasoning:**} \\
\\ 
\texttt{In this hypothetical world, we are given the following structural constraints:} \\
\texttt{1.  \$Y\_\{upt\}\$ has a direct effect on \$Muvq\$.} \\
\texttt{2.  \$Z\_\{upj\}\$ has a direct effect on \$Muvq\$.} \\
\texttt{3.  The statement "We know that \$Y\_\{upt\}\$ or \$Z\_\{upj\}\$ causes \$Muvq\$" implies that the occurrence of \$Muvq\$ is contingent upon the activation of at least one of these two causes. In causal terms, this suggests a disjunctive causal mechanism (an OR gate) where \$Muvq = f(Y\_\{upt\}, Z\_\{upj\})\$ such that if either input is active (assuming "causes" implies the presence of the factor triggers the outcome), the outcome occurs.} \\
\\ 
\texttt{The scenario presents an individual who is **not** \$Z\_\{upj\}\$.} \\
\texttt{*   If the individual is not \$Z\_\{upj\}\$, the causal pathway through \$Z\_\{upj\}\$ is inactive.} \\
\texttt{*   The question asks if the individual would be \$Muvq\$ if \$Y\_\{upt\}\$ were present versus if \$Y\_\{upt\}\$ were absent (not \$Y\_\{upt\}\$).} \\
\\ 
\texttt{Let's analyze the two states for \$Y\_\{upt\}\$ given \$Z\_\{upj\}\$ is false:} \\
\texttt{*   **Case A (\$Y\_\{upt\}\$ is true):** Since \$Y} \\
            }
        }
    }
    & & \\ \\

    \theutterance \stepcounter{utterance}  
    & & & \multicolumn{2}{p{0.3\linewidth}}{
        \cellcolor[rgb]{0.9,0.9,0.9}{
            \makecell[{{p{\linewidth}}}]{
                \texttt{\tiny{[GM$|$GM]}}
                \texttt{No} \\
            }
        }
    }
    & & \\ \\

    \theutterance \stepcounter{utterance}  
    & & & \multicolumn{2}{p{0.3\linewidth}}{
        \cellcolor[rgb]{0.9,0.9,0.9}{
            \makecell[{{p{\linewidth}}}]{
                \texttt{\tiny{[GM$|$GM]}}
                \texttt{yes} \\
            }
        }
    }
    & & \\ \\

    \theutterance \stepcounter{utterance}  
    & & & \multicolumn{2}{p{0.3\linewidth}}{
        \cellcolor[rgb]{0.9,0.9,0.9}{
            \makecell[{{p{\linewidth}}}]{
                \texttt{\tiny{[GM$|$GM]}}
                \texttt{game\_result = LOSE} \\
            }
        }
    }
    & & \\ \\

\end{supertabular}
}

\end{document}
